\section{Fiat Misinterpretations and MMT Fallacies}
\subsection{Structural Misunderstandings}
Conventional fiat practice often inherits pre-1971 gold-standard intuitions, treats sovereign finance as household budgeting, and lacks a unified transparency protocol. These gaps aggravate the GSSH feedback lock by masking true fiscal space and transmission quality.

\subsection{MMT: Descriptive Strength, Policy Exposure}
Modern Monetary Theory accurately notes that taxes stabilize value, bond issuance is an asset swap, and currency-issuing governments avoid involuntary default. Policy exposure arises from reliance on inflation as the sole constraint, limited treatment of cross-border flows, and absence of guardrails against political misuse.

\subsection{ITIF Response}
The International Transparency-Indexed Fiat Framework (ITIF) retains valid descriptive elements while replacing discretionary and opaque components with measurable structures (T-value), transparency-tagged sovereign instruments, and governance that bounds political exposure without creating centralized control.

\subsection{Comparison Table}
\begin{table}[H]
  \centering
  \begin{tabular}{p{0.28\linewidth}p{0.3\linewidth}p{0.3\linewidth}}
    \toprule
    Category & MMT & ITIF \\
    \midrule
    Role of Tax & Maintains nominal value & Value maintenance plus transparency enforcement \\
    Constraint & Inflation only & T-value structural constraint \\
    Bonds & Asset swap & Transparency-tagged sovereign instruments \\
    Political Exposure & Very high & Bounded by CLEAR metrics and public bands \\
    International Applicability & Weak & Strong (IMF/BIS compatible) \\
    Systemic Risk Handling & Assumes state solvency & Multi-factor sovereign integrity model \\
    \bottomrule
  \end{tabular}
  \caption{MMT vs. ITIF}
\end{table}
